\documentclass[11pt]{article}
\usepackage[a4paper,margin=28mm]{geometry}
\usepackage{amsmath,amssymb,amsthm}
\usepackage{graphicx}
\usepackage{microtype}
\usepackage[hidelinks]{hyperref}

\newtheorem{theorem}{Theorem}
\newtheorem{lemma}{Lemma}
\newcommand{\HH}{\mathbb H}
\newcommand{\RR}{\mathbb R}

\title{A Euclidean-stratification upper bound for\\Heisenberg-cylinder discrepancy}
\author{Run \texttt{fa\_banach\_001}, lane 9\\\texttt{agent\_lane\_09}}
\date{11 August 2026}

\begin{document}
\maketitle

\begin{center}
\fbox{\parbox{0.91\textwidth}{\textbf{Status.}
Candidate partial result, likely valid.  The theorem improves the known upper
exponent in a natural Lebesgue-measure subcase; it does not determine the
sharp exponent.}}
\end{center}

\section{Source question}

Brandolini, Monguzzi, and Monti~\cite{BMM} consider
\(\HH^n=\mathbb C^n\times\RR\), with group law
\[
 (z,t)\circ(w,s)=
 \left(z+w,t+s+\tfrac12\operatorname{Im}(z\cdot\overline w)\right),
\]
homogeneous dimension \(Q=2n+2\), and cylinders
\[
 B_\rho=\{(z,t): |z|\leq \rho,\ |t|\leq\rho^2\}.
\]
For a probability measure \(\mu\) and points \(p_1,\ldots,p_N\), they define
\[
 D_N(g;\rho)=\sum_{j=1}^N\mathbf 1_{g\circ B_\rho}(p_j)
       -N\mu(g\circ B_\rho).
\]
Their Theorem~2 proves
\[
 \int_0^1\int_{\HH^n}|D_N(g;\rho)|^2\,dg\,d\rho
 \gtrsim N^{1-1/(Q-2)}
\]
for upper Ahlfors-regular \(\mu\).  They ask whether this is sharp and record,
for suitable normalized Lebesgue measures, the general metric-space upper
bound \(N^{1-1/Q}\).

Figure~\ref{fig:question} is the question on page 3 of the source paper.

\clearpage
\begin{figure}[!ht]
\centering
\includegraphics[height=0.82\textheight]{figures/open_problem_crop.png}
\caption{The sharpness question in arXiv:2601.15850, page 3.}
\label{fig:question}
\end{figure}
\clearpage

\section{Proof intuition}

The source upper bound partitions the support using the intrinsic Heisenberg
metric, whose dimension is \(Q\).  For Lebesgue measure, however, the ambient
Euclidean coordinate dimension is only
\[
 d=2n+1=Q-1.
\]
Partition a cube into equal-volume cells of Euclidean diameter
\(h\asymp N^{-1/d}\), and choose one random point in each cell.  For a fixed
test cylinder, cells entirely inside or outside make no random error; only
cells crossing its boundary contribute variance.  A relevant Heisenberg
translate is an affine shear of a Euclidean ball times an interval.  On the
bounded parameter region that can interact with the cube, these shears have
uniformly bounded condition number, and their boundary layers have Euclidean
volume \(O(h)\).  Thus only \(O(Nh)\) cells contribute.  Averaging over all
translations and radii and then choosing one realization gives the result.

The remaining gap is real: this argument sees ordinary Euclidean codimension
one, not the effective horizontal dimension \(Q-2\) suggested by the lower
bound.

\section{The partial theorem}

Let \(E=[-1/2,1/2]^d\subset\RR^d=\HH^n\) and let \(\mu\) be Lebesgue
measure restricted to \(E\), so \(\mu(E)=1\).

\begin{theorem}\label{thm:main}
For every \(N\geq2\), there exist points \(p_1,\ldots,p_N\in E\) such that
\[
 \int_0^1\int_{\HH^n}|D_N(g;\rho)|^2\,dg\,d\rho
 \leq C_{n,E}N^{1-1/(Q-1)}.
\]
Consequently, in this Lebesgue-measure subcase,
\[
 c_{n,E}N^{1-1/(Q-2)}
 \leq
 \inf_{p_1,\ldots,p_N\in E}
 \int_0^1\int_{\HH^n}|D_N(g;\rho)|^2\,dg\,d\rho
 \leq
 C_{n,E}N^{1-1/(Q-1)}.
\]
In particular, the upper exponent \(1-1/Q\) recorded in the source improves
to \(1-1/(Q-1)\) for this subcase.
\end{theorem}

We use two elementary lemmas.

\begin{lemma}[Equal-volume Euclidean cells]\label{lem:partition}
For every \(d\geq1\) and \(N\geq1\), the unit cube admits a measurable
partition \(E=E_1\sqcup\cdots\sqcup E_N\) such that
\[
 |E_j|=N^{-1},\qquad
 \operatorname{diam}_{\mathrm{Euc}}(E_j)\leq C_dN^{-1/d}.
\]
\end{lemma}

\begin{proof}
This follows by induction on \(d\).  For \(d=1\), use \(N\) equal intervals.
For the inductive step, choose \(m\asymp N^{1/d}\) and write
\(N=N_1+\cdots+N_m\), where all positive integers \(N_j\) are comparable to
\(N/m\).  Divide the first coordinate into slabs of widths \(N_j/N\).  In
the \(j\)-th slab, apply the \((d-1)\)-dimensional construction to its unit
cross-section with \(N_j\) cells.  Every resulting box has volume \(1/N\),
first-coordinate width \(O(N^{-1/d})\), and transverse diameter
\(O(N_j^{-1/(d-1)})=O(N^{-1/d})\).
\end{proof}

\begin{lemma}[Uniform Euclidean boundary layers]\label{lem:layer}
There is a compact set \(K\subset\HH^n\) and a constant \(C\) such that the
following holds for \(0<h\leq1\), \(0\leq\rho\leq1\), and \(g\in\HH^n\).
If the Euclidean \(2h\)-neighborhood of
\(\partial(g\circ B_\rho)\) meets \(E\), then \(g\in K\), and
\[
 \left|E\cap\{x:\operatorname{dist}_{\mathrm{Euc}}
       (x,\partial(g\circ B_\rho))\leq2h\}\right|\leq Ch.
\]
\end{lemma}

\begin{proof}
Write \(g=(z,t)\).  In Euclidean coordinates, left translation is the affine
map
\[
 T_g(w,s)=
 \left(z+w,t+s+\tfrac12\operatorname{Im}(z\cdot\overline w)\right).
\]
It has determinant one.  If a translated cylinder boundary comes within
distance two of the fixed cube, then \(|z|\) is bounded because its horizontal
projection is the ball \(z+\{|w|\leq\rho\}\), with \(\rho\leq1\).  The central
coordinate \(t\) is then bounded as well, since \(|s|\leq\rho^2\) and the shear
term is bounded.  This gives a fixed compact parameter set \(K\).  For
\(g\in K\), the linear parts of \(T_g\) and \(T_g^{-1}\) have uniformly
bounded operator norms.

Before applying \(T_g\), the set \(B_\rho\) is the Euclidean product of a
\(2n\)-ball of radius \(\rho\) and an interval of length \(2\rho^2\).  If
\(\rho\geq h\), its two cap layers and its lateral layer have total volume at
most \(C_nh\), including the edge neighborhoods.  If \(\rho<h\), the whole
\(O(h)\)-neighborhood has volume
\[
 O\big((\rho+h)^{2n}(\rho^2+h)\big)=O(h^{2n+1})=O(h).
\]
The uniformly bounded affine maps \(T_g\) preserve this estimate up to a
constant.  Intersecting with \(E\) can only decrease volume.
\end{proof}

\begin{proof}[Proof of Theorem~\ref{thm:main}]
Take the partition from Lemma~\ref{lem:partition}, and put
\(h=C_dN^{-1/d}\).  The finitely many \(N\) for which \(h>1\) are absorbed by
enlarging the final constant, so assume \(h\leq1\).  Independently in each
\(E_j\), choose \(X_j\) according
to normalized Lebesgue measure on that cell.  For a fixed measurable test set
\(R=g\circ B_\rho\), let
\[
 p_j=N|E_j\cap R|,
 \qquad
 Y_j=\mathbf 1_R(X_j)-p_j.
\]
Then \(\mathbb E Y_j=0\), the variables are independent, and
\[
 \mathbb E|D_N(g;\rho)|^2
 =\sum_{j=1}^N p_j(1-p_j).
\tag{1}
\]
If \(p_j\) is zero or one, its summand vanishes.  Otherwise the cell meets
both \(R\) and its complement.  Since its diameter is at most \(h\), the
whole cell lies in the Euclidean \(2h\)-neighborhood of \(\partial R\).
Lemma~\ref{lem:layer} and the equal cell volumes therefore give
\[
 \#\{j:0<p_j<1\}\,N^{-1}\leq Ch,
\]
and hence, by (1),
\[
 \mathbb E|D_N(g;\rho)|^2\leq CNh.
\tag{2}
\]
Moreover, a nonzero summand in (1) forces the boundary layer to meet \(E\),
so Lemma~\ref{lem:layer} restricts \(g\) to the fixed compact set \(K\).
Tonelli's theorem and (2) yield
\[
 \mathbb E\int_0^1\int_{\HH^n}|D_N(g;\rho)|^2\,dg\,d\rho
 \leq C|K|Nh
 \leq C'N^{1-1/d}.
\]
Some realization of \(X_1,\ldots,X_N\) is no larger than this expectation.
Since \(d=Q-1\), this is the asserted upper bound.

Finally, Lebesgue measure on \(E\) is upper Ahlfors regular for the Kor\'anyi
metric because a Kor\'anyi ball of radius \(r\) has Lebesgue volume
\(c_nr^Q\).  The source's Theorem~2 therefore supplies the displayed lower
bound.
\end{proof}

\section{Verification, limitations, and novelty status}

The proof has no numerical dependency.  It explicitly handles the two
potential uniformity issues: small radii in Lemma~\ref{lem:layer}, and the
unbounded translation domain via the compact active set \(K\).

The result is partial in two ways.  First, it treats normalized Lebesgue
measure on a cube (the same natural type of measure used for the source upper
bound), not every upper Ahlfors-regular probability measure.  Second, it leaves
the exponent gap
\[
 1-\frac1{Q-2}
 \quad\text{versus}\quad
 1-\frac1{Q-1}.
\]
An anisotropic cell count cannot close this gap because the sheared flat caps
vary by first order in both horizontal and central cell widths.

A bounded search covered the exact title and sharpness language, the source's
references on metric discrepancy and halfspaces, the phrases ``Heisenberg
cylinder discrepancy'' and ``quadratic discrepancy Heisenberg sharp'', and
the June 2026 paper on multidimensional nonsmooth convex bodies
\cite{BBM}.  No exact statement of the \(Q-1\) upper bound for this Heisenberg
family was found through 11 August 2026.  The method is nevertheless close to
the standard stratified-sampling argument of Brandolini et al.~\cite{BCCGT},
so novelty confidence is moderate rather than high.

\paragraph{Human-review recommendation.}
Review the uniform affine-shear boundary-layer estimate and confirm that the
source's intended comparison includes the Lebesgue cube subcase.  If valid,
the packet should be retained as a substantial partial upper-bound
improvement, not as a resolution of the conjectured sharp exponent.

\begin{thebibliography}{9}
\bibitem{BMM}
L. Brandolini, A. Monguzzi, and M. Monti,
\emph{Quadratic discrepancy estimates for probability measures on the
Heisenberg group}, arXiv:2601.15850 (2026).

\bibitem{BCCGT}
L. Brandolini, W. W. L. Chen, L. Colzani, G. Gigante, and G. Travaglini,
\emph{Discrepancy and numerical integration on metric measure spaces},
J. Geom. Anal. 29 (2019), 328--369.

\bibitem{BBM}
R. Bramati, L. Brandolini, and A. Monguzzi,
\emph{Discrepancy estimates for multi-dimensional non-smooth convex bodies:
a case study}, arXiv:2606.07028 (2026).
\end{thebibliography}

\end{document}
